\subsection{Digital platforms at scale: attention concentration, portfolio ceilings, and divergent horizons}\label{sec:attraction}

\noindent When partner search moves from long-acquaintance juries to digital platforms, millions of staged profiles leave
an empirical trace of what single-context evaluation amplifies and what life-course intentions it separates. Across two
historical platform exports totaling \PboCombinedRecords{} canonical profile records---\PboHonEligible{} eligible profiles
in the HON facial-rating archive (from \PboHonRows{} canonical accounts and \PboHonFiles{} inventoried images) and
\PboPovEligible{} opposite-sex-seeking adult profiles in the POV dating archive (from \PboPovRows{} unique accounts)---we
measure three phenomena predicted by the two-way jury and episode models (Section~\ref{sec:platformmethods}).

\subsubsection{Beauty age gradients, \PboHonAttentionRatio$\times$ voter skew, and the portfolio rank ceiling}
In the HON archive, within-gender beauty percentiles diverge sharply across recorded adult age while retaining wide
within-age dispersion (Fig.~\ref{fig:pbo-hon-arch}A). At ages 25--29, the mean within-gender beauty percentile is
\PboBeautyYoungW{} for women's profiles and \PboBeautyYoungM{} for men's; by ages 45--49 it falls to \PboBeautyOlderW{} for
women and \PboBeautyOlderM{} for men. Over the same span, the share of profiles reaching the top tenth of their own
gender's rating distribution drops sixteen-fold among women (\PboBeautyTopYoungW{} to \PboBeautyTopOlderW) and four-fold
among men (\PboBeautyTopYoungM{} to \PboBeautyTopOlderM; Fig.~\ref{fig:pbo-hon-arch}B). Because the supplied 0--10 beauty
score tracks the profile's positive-vote share almost deterministically ($r=\PboBeautyRW$ in women, $r=\PboBeautyRM$ in
men), these ranks summarize where profiles stood in voter consensus rather than biological fecundity directly; cross-sectional
platform entry, image curation (\PboImageShareW{} of eligible women and \PboImageShareM{} of men have saved image files),
and unrecorded voter demographics enter the gradient (Fig.~\ref{fig:pbo-ratings}; Table~\ref{tab:pbo-hon-sensitivity}).

Attention on the platform is even more unequal than consensus rank (Fig.~\ref{fig:pbo-hon-arch}C). Eligible women's
profiles receive a median of \PboHonVotesMedW{} votes against \PboHonVotesMedM{} for men---an \PboHonAttentionRatio-fold
asymmetry in evaluative volume. Positive evaluations concentrate sharply in both sexes, and more steeply among men: the
Gini coefficient of received likes is \PboHonLikesGiniW{} among women (vote Gini \PboHonVotesGiniW) and \PboHonLikesGiniM{}
among men (vote Gini \PboHonVotesGiniM), with the top tenth of male profiles capturing \PboHonTopDecileLikesM{} of all
positive votes cast on men (against \PboHonTopDecileLikesW{} among women). Adding more photographs to a profile tests
whether expanding a single-mode visual portfolio raises consensus standing (Fig.~\ref{fig:pbo-hon-arch}D). Moving from one
inventoried photograph to five or more increases median voter traffic by roughly half in both women ($1.51\times$) and men
($1.50\times$), yet mean within-gender beauty percentile stays flat or drifts slightly downward ($33.6$ to $28.8$ in women;
$48.9$ to $47.7$ in men). Consistent with Proposition~\ref{prop:twoway}, extra samples from the same staged photographic
channel buy exposure ($m$), not higher underlying standing ($T_i$).

\begin{figure*}[t]
\centering
\includegraphics[width=\textwidth]{performance-fig-hon-architecture.pdf}
\caption{\textbf{Rating consensus, attention inequality, and portfolio ceilings across \PboHonEligible{} HON profiles.}
Eligible adult profiles aged 18--79 with unconflicted metadata and at least 20 recorded votes. \textbf{(A)}~Mean
within-gender beauty percentile by exact age, with interquartile ranges shaded and the 50th-percentile reference dashed.
\textbf{(B)}~Share of profiles at or above the 90th percentile of their gender's full adult distribution (dashed line:
10\%). \textbf{(C)}~Lorenz concentration curves for received votes (solid) and positive likes (dashed), showing an
\PboHonAttentionRatio-fold female-to-male median vote ratio and steeper top-end concentration among men (top decile receives
\PboHonTopDecileLikesM{} of male likes). \textbf{(D)}~Mean within-gender beauty percentile (solid, left axis) and median
votes relative to single-photo profiles (dashed, right axis) by inventoried photo count: extra photos expand voter attention
by $\sim\!50\%$ without raising consensus rank. Age cells below \PboMinCell{} profiles are omitted.}
\label{fig:pbo-hon-arch}
\end{figure*}

\subsubsection{Canonical facial attribution, super-resolved composites, aging trajectories, and seven-region invariance}\label{sec:faceattr}
What anatomical cues do stranger crowds actually read when a profile photograph is evaluated at a glance, and how much of
voter consensus resides in the face itself rather than framing, metadata, or unmeasured context? To answer both questions
without ever exposing an identifiable individual, we registered \PboFaceTotalN{} stratified primary portraits ($1200\times
1200$ source imagery across sex, five adult age bands, rating deciles, and \PboFaceWorldRegionsN{} world regions) onto a
canonical $106$-landmark, \PboFaceTrisN-triangle piecewise-affine Delaunay atlas~\cite{deng2019arcface} in linear sRGB at
$384\times 384$ studio resolution, enforcing a strict privacy floor ($N\ge 75$ aligned individuals per composite). Combining
three pretrained visual backbones (`DINOv3` spatial tokens~\cite{oquab2024dinov2}, `ConvNeXt V2`~\cite{woo2023convnextv2},
and `AuraFace` `Glint360K-R100`/`SFace`~\cite{deng2019arcface}) with \PboFaceTrisN-triangle Jacobian morphometrics and a
joint $12$-zone interventional Shapley coalition decomposition~\cite{lundberg2017shap} fitted strictly on training profiles
(centered on the neutral 42nd--58th percentile cohort; Section~\ref{sec:platformmethods}), we project consensus beauty into
within-gender percentile drivers ($+\Delta$~pct) and rejectors ($-\Delta$~pct).

First, on the strict profile-grouped $10{,}000$-profile ($20{,}000$-image) benchmark using zero metadata (no age, sex,
geography, or background framing), pure aligned facial morphology and surface appearance achieve held-out Spearman
$\rho=\PboFacePureSpearman$ ($\rho^2=\PboFacePureRankVar$, within-gender percentile MAE \PboFacePurePctMae;
Fig.~\ref{fig:pbo-face-drivers}D and Table~\ref{tab:pbo-ai-bench}). Adding profile metadata and multi-photo aggregation
raises single-channel consensus prediction to $\rho=\PboBenchFinalR$ ($43.8\%$), leaving \PboBenchUnexplainedPct{} of
human-jury rank variance unexplained even at maximum visual resolution---a direct empirical confirmation of
Proposition~\ref{prop:twoway}'s single-context ceiling. Across natural-age held-out evaluation cohorts (\PboFaceEvalNatW{}
women, \PboFaceEvalNatM{} men; live multi-backbone spatial head held-out $\rho=\PboFaceLiveSpearmanW$ in women and
$\rho=\PboFaceLiveSpearmanM$ in men), positive beauty drivers in the top 15\% ($\ge 85\text{th}$ percentile) and negative
rejectors in the bottom 15\% ($\le 15\text{th}$ percentile) localize sharply rather than uniformly across the face
(Fig.~\ref{fig:pbo-face-drivers}A--C; Table~\ref{tab:pbo-face-drivers}). In women, the largest top-minus-bottom leverage sits
in \PboFaceTopZoneNameW{} (top-15\% driver $\PboFaceTopZoneDrvW$ percentile points, bottom-15\% rejector
$\PboFaceTopZoneRejW$, net spread $\PboFaceTopZoneSpreadW$~points), closely followed by eyebrow arch, vermilion lip fullness,
infraorbital tear-trough smoothness, malar cheekbone contour, and jawline definition. In men, leverage is led by
\PboFaceTopZoneNameM{} (driver $\PboFaceTopZoneDrvM$, rejector $\PboFaceTopZoneRejM$, spread $\PboFaceTopZoneSpreadM$~points),
alongside mandibular/chin contour, periorbital aperture, brow structure, and low bilateral asymmetry plus canonical
\PboFaceTrisN-triangle geometric averageness~\cite{langlois1990,rhodes2006,little2011}. Notably, in both sexes bottom-15\%
rejector penalties exceed top-15\% driver gains by roughly twofold across every zone, mirroring the lower-tail floor and
koinophilic averageness ceiling seen for stature (Section~\ref{sec:resheight}; Corollary~\ref{cor:floor}).

\begin{figure*}[!t]
\centering
\includegraphics[width=\textwidth]{performance-fig-face-drivers.pdf}
\caption{\textbf{Super-resolved canonical facial composites, $12$-zone Shapley beauty drivers and rejectors, and the single-context visual ceiling.}
Piecewise-affine $106$-landmark, \PboFaceTrisN-triangle Delaunay registration across natural-age held-out evaluation cohorts
(\PboFaceEvalNatW{} women, \PboFaceEvalNatM{} men; every composite averages $N\ge 190$ aligned portraits). \textbf{(A,~B)}~Linear-sRGB
super-resolved studio composites for the bottom 15\%, median (42nd--58th), and top 15\% within-gender beauty tiers, paired with
signed percentile rejector maps ($-\Delta$~pct, sapphire), driver maps ($+\Delta$~pct, coral-gold), and \PboFaceTrisN-triangle
morphometric expansion/contraction contrasts. \textbf{(C)}~Joint $12$-zone interventional Shapley contributions (mean $\Delta$
within-gender percentile points with 600-draw bootstrap 95\% CIs) separating positive drivers in the top 15\% (solid) from
negative rejectors in the bottom 15\% (hatched). \textbf{(D)}~Explained human-jury rank variance ($\rho^2$) on the
$10{,}000$-profile ($20{,}000$-image) held-out benchmark from raw UV crops (H1) to pure aligned facial anatomy (zero metadata:
$\rho=\PboFacePureSpearman$, \PboFacePureRankVar) and the full multimodal stack (H20: $43.8\%$), versus the
\PboBenchUnexplainedPct{} single-context gap.}
\label{fig:pbo-face-drivers}
\end{figure*}

\begin{table*}[!t]
\caption{Anatomical beauty drivers ($+\Delta$ within-gender percentile points in the top 15\%) and rejectors ($-\Delta$ points
in the bottom 15\%) across $12$ canonical Delaunay facial zones (natural-age held-out evaluation cohorts; 600-draw bootstrap
95\% confidence intervals in brackets).}\label{tab:pbo-face-drivers}
\scriptsize
\centering
\begin{tabularx}{\textwidth}{@{}Lcccccc@{}}
\toprule
\tabinput{performance-face-drivers-table.tex}
\bottomrule
\end{tabularx}
\tablenote{Zones are defined on the symmetric $106$-landmark canonical ArcFace/Delaunay template and evaluated by joint
$12$-zone interventional Shapley coalition attribution centered on the neutral 42nd--58th percentile training cohort.}
\end{table*}

Second, stratifying non-training portraits across five adult age buckets (`18--24`, `25--34`, `35--44`, `45--54`, and `55+`;
$n\ge 220$ aligned faces per sex-by-age cell) isolates the physical basis of the platform age gradient
(Fig.~\ref{fig:pbo-face-age}). Across life-course composites, observed voter percentile declines from
$\PboFaceAgeObsYoungW\text{th}$ at ages 18--24 to $\PboFaceAgeObsOldW\text{th}$ at $55+$ in women (pure-face model prediction:
$\PboFaceAgePredYoungW\text{th}$ to $\PboFaceAgePredOldW\text{th}$) and from $\PboFaceAgeObsYoungM\text{th}$ to
$\PboFaceAgeObsOldM\text{th}$ in men (pure-face prediction: $\PboFaceAgePredYoungM\text{th}$ to
$\PboFaceAgePredOldM\text{th}$; Fig.~\ref{fig:pbo-face-age}B). Signed attribution maps show warm periorbital, malar, and
vermilion drivers dominating young-adult portraits before transitioning progressively after age 35 toward cool infraorbital
tear-trough, nasolabial, and lower-facial soft-tissue rejectors, with steeper periorbital and midface rank penalties in women
than in men (Fig.~\ref{fig:pbo-face-age}A,C).

\begin{figure*}[!t]
\centering
\includegraphics[width=\textwidth]{performance-fig-face-age.pdf}
\caption{\textbf{Super-resolved facial aging composites and signed beauty-rank attribution across five adult age cohorts.}
Non-training evaluation portraits stratified across ages 18--24 through $55+$ ($n\ge 220$ aligned individuals per cell).
\textbf{(A)}~Linear-sRGB canonical composites (top rows per sex) and signed within-gender beauty percentile attribution maps
(bottom rows per sex), showing the spatial transition from young-adult periorbital, malar, and lip drivers (coral-gold) to
midlife infraorbital, nasolabial, and lower-facial rejectors (sapphire). \textbf{(B)}~Mean observed human-voter beauty
percentile (solid) versus metadata-blind pure-face predicted percentile (dashed) across age groups. \textbf{(C)}~Net anatomical
zone drift from ages 18--24 to $55+$ ($\Delta$ within-gender percentile contribution).}
\label{fig:pbo-face-age}
\end{figure*}

Third, computing super-resolved canonical portraits and within-region beauty driver maps (top half minus bottom half of local
voter percentiles) across all \PboFaceWorldRegionsN{} world regions---United States, South America, Mexico/Central America \&
Caribbean, Canada, Asia-Pacific \& Oceania, Europe, and the Middle East \& Africa---tests whether facial evaluation criteria
fragment across geography (Fig.~\ref{fig:pbo-face-world}). Whereas regional upload selectivity shifts raw voter means
(Fig.~\ref{fig:pbo-face-world}C), the spatial anatomy of within-region beauty drivers is strikingly invariant across the globe:
pairwise correlations across the \PboFaceWorldRegionsN{} regional driver maps average $r=\PboFaceWorldCorrW$ in women
($\min r=\PboFaceWorldMinCorrW$) and $r=\PboFaceWorldCorrM$ in men ($\min r=\PboFaceWorldMinCorrM$;
Fig.~\ref{fig:pbo-face-world}B). Single-context visual crowds everywhere read essentially the same morphological cues of youth,
ocular-lip definition, and symmetry---and everywhere stop at the same single-channel horizon when forecasting long-term conduct.

\begin{figure*}[!t]
\centering
\includegraphics[width=\textwidth]{performance-fig-face-world.pdf}
\caption{\textbf{Global regional canonical composites and cross-regional invariance of facial beauty drivers across \PboFaceWorldRegionsN{} world regions.}
Natural-age non-training portraits across United States, South America, Mexico/Central America \& Caribbean, Canada,
Asia-Pacific \& Oceania, Europe, and Middle East \& Africa ($n\ge 75$ aligned portraits per regional composite).
\textbf{(A)}~Super-resolved linear-sRGB regional portraits and within-region beauty driver maps (above-median minus below-median
within-gender attribution). \textbf{(B)}~Pairwise spatial correlation matrix ($r$) across regional beauty driver maps (lower
triangle: women, mean $r=\PboFaceWorldCorrW$; upper triangle: men, mean $r=\PboFaceWorldCorrM$). \textbf{(C)}~Regional mean
human-voter percentile (filled) versus metadata-blind pure-face prediction (open), separating upload self-selection from
anatomical scoring.}
\label{fig:pbo-face-world}
\end{figure*}

\subsubsection{Childbearing intentions reverse with age while relationship commitment does not}
The POV dating archive separates two horizons that fertility discussions often conflate: wanting future children and wanting
a committed partnership (Fig.~\ref{fig:pbo-pov-life}). Among childless profiles seeking the other recorded gender,
\PboPovWantYoungW{} of women and \PboPovWantYoungM{} of men aged 25--29 select ``Want children'' (women minus men:
\PboPovGapYoung{} percentage points). By ages 45--49 the shares cross to \PboPovWantOlderW{} of women and
\PboPovWantOlderM{} of men (\PboPovGapOlder{} points; Fig.~\ref{fig:pbo-pov-life}A). The reversal holds under exact-age
standardization, among profiles marked single, across all seeking categories, and after excluding nonanswers or
``Undecided/Open'' responses (\PboPovOpenYoungW{} and \PboPovOpenYoungM{} at 25--29; \PboPovOpenOlderW{} and
\PboPovOpenOlderM{} at 45--49; Tables~\ref{tab:pbo-pov-responses} and~\ref{tab:pbo-pov-sensitivity};
Fig.~\ref{fig:pbo-pov}).

Relationship commitment moves in the opposite way: it never reverses (Fig.~\ref{fig:pbo-pov-life}B). Women exceed men in
seeking marriage or a serious relationship in every age window across the adult life course---from \PboPovSeriousYoungW{}
versus \PboPovSeriousYoungM{} at ages 25--29, to \PboPovSeriousMidW{} versus \PboPovSeriousMidM{} at 40--49, to
\PboPovSeriousSenW{} versus \PboPovSeriousSenM{} at 60--69. Women on the platform also report more prior union durability
at midlife (\PboPovTenYrMidW{} of women vs.\ \PboPovTenYrMidM{} of men aged 40--49 report a longest relationship exceeding
ten years, as the male-to-female profile ratio narrows from $2.95{:}1$ at 25--29 to $0.98{:}1$ at 60--69;
Fig.~\ref{fig:pbo-pov-life}C). Within prime reproductive and career-building years, self-reported ambition and schooling
trace postponement rather than rejection of parenthood (Fig.~\ref{fig:pbo-pov-life}D,E). Among women aged 25--49, moving
from ``Not Ambitious'' to ``Very Ambitious'' is associated with lower current motherhood (\PboPovAmbParentLowW{} to
\PboPovAmbParentHighW) alongside \emph{higher} stated desire for children (\PboPovAmbWantLowW{} to \PboPovAmbWantHighW);
across completed education at ages 30--49, current parenthood falls steeply from high school to doctoral degrees in both
sexes while childless women's desire for children rises through master's-level schooling.

Crucially, this decoupling of biological fertility timing from long-term partner commitment is not an American idiosyncrasy.
Across all five countries with large POV sub-cohorts---the United States, United Kingdom, Canada, Brazil, and
Australia---childless ``Want children'' flips from positive (women above men by $+4.2$ to $+19.7$ points at ages 25--29) to
negative in all five countries at ages 45--49 ($-3.7$ to $-13.3$ points), whereas serious relationship intent remains
positive in all five countries at both 25--29 ($+8.8$ to $+15.1$ points) and 45--49 ($+7.1$ to $+16.4$ points;
Fig.~\ref{fig:pbo-pov-life}F and Table~\ref{tab:pbo-pov-geo}). Midlife women do not stop seeking committed partners; they
adjust fertility intentions to a finite reproductive window while men's stated fertility horizon extends later.

\begin{figure*}[t]
\centering
\includegraphics[width=\textwidth]{performance-fig-pov-lifecourse.pdf}
\caption{\textbf{Divergent fertility and commitment horizons across \PboPovEligible{} POV dating profiles.} Opposite-sex
seeking profiles aged 18--79. \textbf{(A)}~Share selecting ``Want children'' among childless profiles by age band (shading:
25--29 and 45--49). \textbf{(B)}~Share seeking marriage or a committed relationship across six adult age windows: women
exceed men at every age despite the fertility-intention crossover. \textbf{(C)}~Share reporting a prior relationship
lasting over ten years, alongside the narrowing male-to-female platform participation ratio. \textbf{(D)}~Current
parenthood and stated desire for children by self-reported ambition at ages 25--49. \textbf{(E)}~Current parenthood (solid)
and childless ``Want children'' (dashed) across completed education at ages 30--49, showing reproductive postponement rather
than lower family ideals among more educated women. \textbf{(F)}~Replication across all five major country cohorts:
women-minus-men gaps in childless ``Want children'' reverse from ages 25--29 (green circles) to 45--49 (red squares) in 5/5
countries, while serious commitment at 45--49 (violet diamonds) remains positive in 5/5.}
\label{fig:pbo-pov-life}
\end{figure*}

\begin{table*}[t]
\caption{Five-country replication of the childbearing-intention reversal and persistent female commitment advantage in the
POV archive (opposite-sex-seeking profiles). Entries in the last four columns are women-minus-men differences in percentage
points.}\label{tab:pbo-pov-geo}
\footnotesize
\begin{tabularx}{\textwidth}{@{}Lrrrrrr@{}}
\toprule
\tabinput{performance-pov-geography.tex}
\bottomrule
\end{tabularx}
\tablenote{Every cell in the five-country sub-cohorts exceeds \PboMinCell{} profiles. Serious commitment combines marriage
and long-term relationship intent.}
\end{table*}
